% GENERATED FILE. Edit canonical lessons, then run npm run generate:books.
% canonical-source-hash: fnv1a64:4df34f1a59315c36
% canonical-lessons: UR-C38-hath, UR-C38-sar, UR-C38-bal, UR-C38-dant, UR-C38-gala

\chapter{Hand, Head, Hair, Tooth, Throat}
\label{ch:ur-hand-head-hair-tooth-throat}
\hlchaptermodality{38}

\begin{chapteropening}
\textbf{By the end of this chapter:} \emph{I can say a hand, the head, hair, a tooth and the throat.}

Complete the last lesson of chapter 38: I can say a hand, the head, hair, a tooth and the throat.
\end{chapteropening}

\section[\texorpdfstring{\ur{ہاتھ} (hāth)}{ہاتھ (hāth)}]{\ur{ہاتھ} (hāth) — a hand}

\label{lesson:UR-C38-hath}



\begin{quote}
Before the new one: say the Urdu for third, then the Urdu for fourth.
\end{quote}

\subsection*{You'll want to know: \ur{ہاتھ}}
\textbf{\ur{ہاتھ}} — \emph{hāth} — \textquotedblleft{}a hand\textquotedblright{}.

\begin{cousinweb}
\textbf{\ur{ہاتھ}} (\emph{hāth}) is a tadbhava, a word inherited from Sanskrit and reshaped by sound change. It comes from \emph{hasta}, \textquotedblleft{}hand\textquotedblright{} (CDIAL 14024).
\end{cousinweb}

\begin{grammarlens}[title={What you've built}]
One more word for this chapter.
\end{grammarlens}

\subsection*{Your turn}
\begin{itemize}
\raggedright
  \item \emph{Say it:} \emph{hāth}
  \item \emph{Say it:} \emph{hāth}, once more
  \item \emph{Say it:} say \emph{hāth}
  \item \emph{Recall:} say the Urdu for third, then the Urdu for fourth, then say \emph{hāth} again
\end{itemize}

\begin{tcolorbox}[breakable,colback=teal!4,colframe=teal!35!black,title={Before you move on}]
What does \ur{ہاتھ} mean? (\textquotedblleft{}A hand\textquotedblright{}.) Say it once more.
\end{tcolorbox}
\section[\texorpdfstring{\ur{سر} (sar)}{سر (sar)}]{\ur{سر} (sar) — the head}

\label{lesson:UR-C38-sar}



\begin{quote}
Before the new one: say the Urdu for fourth, then the Urdu for a hand.
\end{quote}

\subsection*{You'll want to know: \ur{سر}}
\textbf{\ur{سر}} — \emph{sar} — \textquotedblleft{}the head\textquotedblright{}.

\begin{cousinweb}
\textbf{\ur{سر}} (\emph{sar}) is a loan, usually traced to Persian \emph{sar}, \textquotedblleft{}head\textquotedblright{}. Persian \emph{sar} and Sanskrit \emph{śiras}, the source of Hindi \emph{sir} (CDIAL 12452), are usually taken to be the same ancient word, kept separately by the two branches.
\end{cousinweb}

\begin{grammarlens}[title={What you've built}]
One more word for this chapter.
\end{grammarlens}

\subsection*{Your turn}
\begin{itemize}
\raggedright
  \item \emph{Say it:} \emph{sar}
  \item \emph{Say it:} \emph{sar}, once more
  \item \emph{Say it:} say \emph{sar}
  \item \emph{Recall:} say the Urdu for fourth, then the Urdu for a hand, then say \emph{sar} again
\end{itemize}

\begin{tcolorbox}[breakable,colback=teal!4,colframe=teal!35!black,title={Before you move on}]
What does \ur{سر} mean? (\textquotedblleft{}The head\textquotedblright{}.) Say it once more.
\end{tcolorbox}
\section[\texorpdfstring{\ur{بال} (bāl)}{بال (bāl)}]{\ur{بال} (bāl) — hair}

\label{lesson:UR-C38-bal}



\begin{quote}
Before the new one: say the Urdu for a hand, then the Urdu for the head.
\end{quote}

\subsection*{You'll want to know: \ur{بال}}
\textbf{\ur{بال}} — \emph{bāl} — \textquotedblleft{}hair\textquotedblright{}.

\begin{cousinweb}
\textbf{\ur{بال}} (\emph{bāl}) is a tadbhava, a word inherited from Sanskrit and reshaped by sound change. It comes from \emph{vāla}, \textquotedblleft{}hair of the tail, hair\textquotedblright{} (CDIAL 11572).
\end{cousinweb}

\begin{grammarlens}[title={What you've built}]
One more word for this chapter.
\end{grammarlens}

\subsection*{Your turn}
\begin{itemize}
\raggedright
  \item \emph{Say it:} \emph{bāl}
  \item \emph{Say it:} \emph{bāl}, once more
  \item \emph{Say it:} say \emph{bāl}
  \item \emph{Recall:} say the Urdu for a hand, then the Urdu for the head, then say \emph{bāl} again
\end{itemize}

\begin{tcolorbox}[breakable,colback=teal!4,colframe=teal!35!black,title={Before you move on}]
What does \ur{بال} mean? (\textquotedblleft{}Hair\textquotedblright{}.) Say it once more.
\end{tcolorbox}
\section[\texorpdfstring{\ur{دانت} (dānt)}{دانت (dānt)}]{\ur{دانت} (dānt) — a tooth}

\label{lesson:UR-C38-dant}



\begin{quote}
Before the new one: say the Urdu for the head, then the Urdu for hair.
\end{quote}

\subsection*{You'll want to know: \ur{دانت}}
\textbf{\ur{دانت}} — \emph{dānt} — \textquotedblleft{}a tooth\textquotedblright{}.

\begin{cousinweb}
\textbf{\ur{دانت}} (\emph{dānt}) is a tadbhava, a word inherited from Sanskrit and reshaped by sound change. It comes from \emph{danta}, \textquotedblleft{}tooth\textquotedblright{} (CDIAL 6152). It is the same ancient word as Latin \emph{dēns} and English \emph{tooth}.
\end{cousinweb}

\begin{grammarlens}[title={What you've built}]
One more word for this chapter.
\end{grammarlens}

\subsection*{Your turn}
\begin{itemize}
\raggedright
  \item \emph{Say it:} \emph{dānt}
  \item \emph{Say it:} \emph{dānt}, once more
  \item \emph{Say it:} say \emph{dānt}
  \item \emph{Recall:} say the Urdu for the head, then the Urdu for hair, then say \emph{dānt} again
\end{itemize}

\begin{tcolorbox}[breakable,colback=teal!4,colframe=teal!35!black,title={Before you move on}]
What does \ur{دانت} mean? (\textquotedblleft{}A tooth\textquotedblright{}.) Say it once more.
\end{tcolorbox}
\section[\texorpdfstring{\ur{گلا} (galā)}{گلا (galā)}]{\ur{گلا} (galā) — the throat}

\label{lesson:UR-C38-gala}



\begin{quote}
Before the new one: say the Urdu for hair, then the Urdu for a tooth.
\end{quote}

\subsection*{You'll want to know: \ur{گلا}}
\textbf{\ur{گلا}} — \emph{galā} — \textquotedblleft{}the throat\textquotedblright{}.

\begin{cousinweb}
\textbf{\ur{گلا}} (\emph{galā}) is a tadbhava, a word inherited from Sanskrit and reshaped by sound change. It comes from \emph{gala}, \textquotedblleft{}throat, neck\textquotedblright{} (CDIAL 4070).
\end{cousinweb}

\begin{grammarlens}[title={What you've built}]
One more word for this chapter.
\end{grammarlens}

\subsection*{Your turn}
\begin{itemize}
\raggedright
  \item \emph{Say it:} \emph{galā}
  \item \emph{Say it:} \emph{galā}, once more
  \item \emph{Say it:} say \emph{galā}
  \item \emph{Recall:} say the Urdu for hair, then the Urdu for a tooth, then say \emph{galā} again
\end{itemize}

\begin{tcolorbox}[breakable,colback=teal!4,colframe=teal!35!black,title={Before you move on}]
What does \ur{گلا} mean? (\textquotedblleft{}The throat\textquotedblright{}.) Say it once more.
\end{tcolorbox}
